\section{Acoustic Wave Equation}\label{sec:acoustic-helmholtz}

The propagation of acoustic waves in a homogeneous, inviscid fluid medium is governed by the linear wave equation for the acoustic pressure $p(\mathbf{x}, t)$:

\newglossaryentry{pxt}{
    name={$p(x, t)$},
    description = {Time dependent acoustic pressure}
}
\begin{equation}
    \nabla^2 p(\mathbf{x}, t) - \frac{1}{c^2} \frac{\partial^2 p(\mathbf{x}, t)}{\partial t^2} = 0
    \label{eq:wave_equation}
\end{equation}

\newglossaryentry{x}{
    name={$\mathbf{x}$},
    description = {Position vector (field point)}
}

\newglossaryentry{t}{
    name={$t$},
    description = {Time}
}
\newglossaryentry{c}{
    name={$c$},
    description = {Speed of sound in the medium}
}

where $\mathbf{x}$ is the position vector, $t$ is time, and $c$ is the speed of sound in the medium.

\newglossaryentry{P}{
    name={$P(\mathbf{x})$},
    description = {Spatial component of the acoustic pressure (phasor)}
}

\newglossaryentry{smallomega}{
    name={$\omega$},
    description = {Angular frequency}
}
\newglossaryentry{i}{
    name={$i$},
    description = {Imaginary unit, $\sqrt{-1}$}
}


To derive the Helmholtz equation, we consider a time-harmonic acoustic field. We assume the pressure can be represented as a complex function oscillating with angular frequency $\omega$:

\begin{equation}
    p(\mathbf{x}, t) = P(\mathbf{x}) e^{-i \omega t}
    \label{eq:time_harmonic}
\end{equation}

where $P(\mathbf{x})$ is the spatial component of the pressure (the phasor) and $i = \sqrt{-1}$ is the imaginary unit.

Substituting Equation (\ref{eq:time_harmonic}) into the wave equation (\ref{eq:wave_equation}):

\begin{equation}
    \nabla^2 \left( P(\mathbf{x}) e^{-i \omega t} \right) - \frac{1}{c^2} \frac{\partial^2}{\partial t^2} \left( P(\mathbf{x}) e^{-i \omega t} \right) = 0
\end{equation}

Calculating the time derivatives:
\begin{align}
    \frac{\partial}{\partial t} \left( P(\mathbf{x}) e^{-i \omega t} \right) &= -i \omega P(\mathbf{x}) e^{-i \omega t} \\
    \frac{\partial^2}{\partial t^2} \left( P(\mathbf{x}) e^{-i \omega t} \right) &= (-i \omega)^2 P(\mathbf{x}) e^{-i \omega t} = -\omega^2 P(\mathbf{x}) e^{-i \omega t}
\end{align}

Substituting this back into the wave equation:

\begin{equation}
    e^{-i \omega t} \nabla^2 P(\mathbf{x}) - \frac{1}{c^2} (-\omega^2 P(\mathbf{x}) e^{-i \omega t}) = 0
\end{equation}

Dividing by the common time factor $e^{-i \omega t}$ (which is non-zero) and rearranging the terms:

\begin{equation}
    \nabla^2 P(\mathbf{x}) + \frac{\omega^2}{c^2} P(\mathbf{x}) = 0
\end{equation}

\newglossaryentry{k}{
    name={$k$},
    description = {Wavenumber, defined as $k = \omega / c$}
}

Introducing the wavenumber $k = \omega / c$, we arrive at the homogeneous Helmholtz equation:

\begin{equation}
    \nabla^2 P(\mathbf{x}) + k^2 P(\mathbf{x}) = 0
    \label{eq:helmholtz}
\end{equation}

\newglossaryentry{phi}{
    name={$\phi(\mathbf{x})$},
    description = {Acoustic potential, represents the acoustic pressure in the BEM formulation}
}

Since we use the symbol $\phi$ to represent the acoustic potential in our BEM formulation, we can rewrite the Helmholtz equation as:
\begin{equation}
    \nabla^2 \phi(\mathbf{x}) + k^2 \phi(\mathbf{x}) = 0
    \label{eq:helmholtz_phi}
\end{equation}